\chapter{SONUÇLAR VE GELECEK ÇALIŞMALAR}
\label{chap:sonuc}

Tezde T-tipi dönüştürücünün düşük güç davranışı, nötr anahtar susturma yöntemi, Sin\textsuperscript{3} tabanlı ileri besleme kompanzasyonu ve GT-tipi yeniden yapılandırılabilir enerji depolama mimarisi tek bir araştırma zinciri altında ele alınmıştır. Sonuçlar birbirini tamamlayan kanıt türlerine dayanmaktadır. Düşük güç anahtarlama ve Sin\textsuperscript{3} incelemeleri analitik model, deney özetine kalibre edilmiş anahtarlamalı benzetim, kapalı çevrim S4 benzetimi ve prototip deney kayıtlarına dayanmaktadır. GT-tipi inceleme ise aynı deney düzeneğinin evirici kipindeki deneylerine, şarj zinciri benzetim kayıtlarına ve topoloji tanımına dayanmaktadır. Sonuçlar, dayandıkları kanıt türüne göre doğrulanmış mekanizma, raporlanmış sayısal bulgu, yeniden üretim/uyumluluk sonucu, analitik sonuç ve mimari tanım olarak sınıflandırılmıştır.

\section{Araştırma Sorularının Yanıtları}

Tablo~\ref{tab:arastirma-sorusu-kanit}, Giriş bölümünde tanımlanan dört araştırma sorusunu yanıtlarıyla ve dayandıkları kanıtlarla eşleştirmektedir. ``Sonuç düzeyi'' sütunu yanıtın kapsamını, ``Kanıt türü'' sütunu ise yanıtın deney, kalibre benzetim, analitik model ya da tasarım tanımından hangisine dayandığını göstermektedir.

\begin{longtable}{@{}>{\raggedright\arraybackslash}p{0.07\textwidth}>{\raggedright\arraybackslash}p{0.23\textwidth}>{\raggedright\arraybackslash}p{0.28\textwidth}>{\raggedright\arraybackslash}p{0.28\textwidth}@{}}
  \caption{Araştırma sorusu--kanıt--sonuç matrisi.}
  \label{tab:arastirma-sorusu-kanit}\\
  \toprule
  \textbf{Soru} & \textbf{Sonuç düzeyi} & \textbf{Başlıca dayanak} & \textbf{Kanıt türü} \\
  \midrule
  \endfirsthead
  \caption[]{Araştırma sorusu--kanıt--sonuç matrisi (devam).}\\
  \toprule
  \textbf{Soru} & \textbf{Sonuç düzeyi} & \textbf{Başlıca dayanak} & \textbf{Kanıt türü} \\
  \midrule
  \endhead
  \bottomrule
  \endfoot
  AS1 & Mekanizma düzeyinde yanıtlandı. & Bölüm~\ref{chap:sistem}'deki anahtarlama durumları, akım yolları ve senkron-buck eşdeğeri; Bölüm~\ref{chap:anahtarlama}'daki DCM ve ters akım incelemesi. & Devre analizi (anahtarlama durumları, akım yolları, senkron-buck eşdeğeri) ve kalibre anahtarlamalı benzetimin $L_1$ akım kaydı. \\
  AS2 & Düşük güçte (yaklaşık \SI{2.3}{\ampere}) deney ve benzetimle, orta ve anma akımında S4 benzetimiyle yanıtlandı. & Tablo~\ref{tab:kalibrasyon-karsilastirma} ve Tablo~\ref{tab:deney-dusuk-guc-ozeti}; \SI{10}{\ampere} ve \SI{25}{\ampere} için Tablo~\ref{tab:s4-tarama}. & Deney ve kalibre benzetim (yaklaşık \SI{2.3}{\ampere}), S4 kapalı çevrim benzetimi (\SI{10}{\ampere} ve \SI{25}{\ampere}). \\
  AS3 & Harmonik etkisi düşük güçte deney ve benzetimle, kararlılık ve anma akımında kapalı çevrim izleme analitik model ve benzetimle yanıtlandı. & Tablo~\ref{tab:kalibrasyon-karsilastirma}, Şekil~\ref{fig:sin3-yeniden-tarama}, Tablo~\ref{tab:s4-tarama}, Tablo~\ref{tab:kararlilik-marjlar} ve Tablo~\ref{tab:kararlilik-olcum}. & Deney ve kalibre benzetim (harmonik etki), analitik model (katsayı taraması), S4 kapalı çevrim benzetimi ve analog eşdeğer çevrimin kararlılık analizi (anma akımı ve marjlar). \\
  AS4 & Mimari ve güvenli çalışma koşulları düzeyinde yanıtlandı. & Bölüm~\ref{chap:gttipi}'deki işlev paylaşım yolları, şarj zinciri benzetim kayıtları, çalışma modu ve doğrulama tabloları; evirici kipinde Bölüm~\ref{chap:deney} deneyleri. & Deney (evirici kipi), benzetim kaydı (şarj zinciri) ve tasarım tanımı (çalışma modları, güvenli kip geçişi). \\
\end{longtable}

\textbf{AS1 -- Ters akım mekanizması.} T-tipi faz kolunda sıfır gerilim durumunu oluşturan nötr anahtarın, endüktör akımı sıfıra ulaştıktan sonra iletimde tutulması senkron-buck benzeri çift yönlü bir akım yolu oluşturmaktadır. Akım referansla aynı yönde yararlı güç taşımak yerine ters yönde büyüyebildiği için RMS akım ve iletim kaybı potansiyeli artmaktadır. Bölüm~\ref{chap:sistem}'deki ortak tabanlı akım yolu çizimleri ile Bölüm~\ref{chap:anahtarlama}'daki kip sınırı fiziksel mekanizmayı açıklamaktadır.

\textbf{AS2 -- Nötr anahtar susturmanın verim--harmonik ödünleşimi.} Yaklaşık \SI{2.3}{\ampere} faz akımında susturma, klasik anahtarlamaya göre verimi deneyde \%92,5'ten \%94,2'ye (1,7 puan), deneyin aktif ve reaktif gücüne kalibre edilmiş benzetimde \%91,15'ten \%93,93'e (2,78 puan) yükseltmiştir (Tablo~\ref{tab:kalibrasyon-karsilastirma}). Verim sıralaması iki kaynakta aynıdır; benzetimdeki fazla artış S1 verimindeki $-1{,}35$ puanlık sapmadan gelmektedir. Harmonik bedel iki kaynakta da açıktır: şebeke akımı THD'si deneyde \%15'ten \%47'ye, kalibre benzetimde \%39,5'ten \%56,7'ye çıkmıştır. S4 benzetiminde \SI{10}{\ampere}'de susturmalı köprü çevrimin büyük kısmında süreksiz iletimdedir ve THD \%19,6--22,6 düzeyindedir (Tablo~\ref{tab:s4-tarama}). Ulaşılan sonuç, susturmanın ters akım yolunu keserek düşük güçte verimi artırdığı ve aynı noktada üçüncü harmonik ağırlıklı bir bozulma oluşturduğudur.

\textbf{AS3 -- Sin\textsuperscript{3} kompanzasyonu ve kararlılık.} Susturma ile birlikte kullanılan Sin\textsuperscript{3} terimi, aynı düşük güç noktasında deneyde akım THD'sini \%47'den \%14'e indirmiş ve verimi \%94,5'e taşımıştır. THD klasik anahtarlamanın \%15 düzeyine dönerken verim klasik anahtarlamadan 2,0 puan yüksek kalmıştır; osiloskop kayıtlarında üçüncü harmonik oranı yaklaşık \%47'den yaklaşık \%9'a inmektedir (Alt Bölüm~\ref{sec:thd-kaynaklar}). Deney özetine kalibre edilen S3 modeli deneyin \%14 akım THD'sini kalibre edilmeden \%14,8 ile, verimini 0,25 puan ve AA akımını \%1,4 içinde yeniden üretmektedir (Alt Bölüm~\ref{sec:kalibrasyon}). Analitik taramada $(k_1,k_3)=(0{,}34,\ 0{,}40)$ çifti üçüncü harmonik oranını \%39,96'dan \%0,15'e düşürmüştür; uygun katsayılar çalışma noktasıyla değişir (Alt Bölüm~\ref{sec:sin3-yeniden-tarama}). Anma akımında (\SI{25}{\ampere}) PI ve PIR denetimi doyuma girmeden $-$\%2,1 ve $-$\%1,0 genlik hatasında kalmıştır (Alt Bölüm~\ref{sec:s4-kapali-cevrim}). Kararlılık analizi Sin\textsuperscript{3} teriminin karakteristik denklemde yer almadığını, $L$, $S$ ve $T$'nin $K_{34}$ açık ve kapalıyken aynı kaldığını ve analog eşdeğer çevrimin kritik geçişte $54^{\circ}$ faz marjı ve $+$\SI{35}{\decibel} kazanç marjıyla kararlı olduğunu göstermiştir (Alt Bölüm~\ref{sec:kararlilik-bode}). Anma akımı ve kararlılık sonuçları S4 kapalı çevrim modeli ve kararlılık analiziyle elde edilmiştir.

\textbf{AS4 -- GT-tipi yeniden görevlendirme.} İşlev paylaşım topolojisi; evirici/LCL, iki seviyeli evirici--transformatör ve PFC yollarının aynı yarı iletken tabanının farklı kiplerde kullanılabileceğini devre düzeyinde göstermektedir. Deney düzeneği GT-tipi olarak kurulmuştur; düşük güç deneyleri donanımın evirici kipinde yapılmış, şarj yönünün işlevsel bütünlüğü benzetim kayıtlarıyla gösterilmiştir. RL1 kip geçişinin güvenliği PWM'in durdurulmasına, faz akımlarının sıfıra yaklaşmasına, DA bara dengesine ve röle geri bildirimine birlikte bağlanmıştır.

\section{Özgün Katkıların Değerlendirilmesi}

Alt Bölüm~\ref{sec:giris-katkilar} altındaki Tablo~\ref{tab:katki-yontem-kanit}, üç katkıyı tez başlığının üç bileşeniyle, yöntemle ve kanıtla eşleştirmiştir. Alt kesimler, her katkının tezin sonunda ulaştığı kanıt düzeyini aynı sırayla değerlendirmektedir.

\subsection{Düşük Güç Anahtarlama Katkısı}

T-tipi eviricinin düşük güçteki ters akım sorunu genel bir hafif yük kaybı olarak bırakılmamış; sıfır gerilim durumu, nötr kol ve akım yönü üzerinden topolojiye özgü bir mekanizma olarak açıklanmıştır. Nötr anahtar susturma yaklaşımı, akım sıfıra ulaştıktan sonra ters yönde büyümeyi sürdüren komutu engelleyen işlevsel bir yöntem sunmaktadır. Yöntemin kapı düzeyindeki tanımı modelden doğrulanmış (Alt Bölüm~\ref{sec:model-kapi-mantigi}), verim kazancı deney ve kalibre benzetimde aynı yönde bulunmuştur (AS2). Yöntemin çalışma bölgesi de aynı kanıtla belirlenmiştir: susturma düşük güçte verim kazancı sağlar. \SI{10}{\ampere}'de süreksiz iletim THD'yi \%20 düzeyine taşıdığından orta akımda tamamlayıcı anahtarlama tercih edilir.

\subsection{Sin\textsuperscript{3} İleri Besleme Katkısı}

Sin\textsuperscript{3} terimi bağımsız bir geri besleme kontrolcüsü veya genel amaçlı harmonik süzgeç olarak değil, nötr susturmanın ürettiği üçüncü harmonik cezasını azaltan bir modülasyon ileri beslemesi olarak konumlandırılmıştır. Tez, $\sin^3(\theta)$ ile $\sin(3\theta)$ katsayılarının aynı anlamı taşımadığını göstermiş; katsayıların tek bir cebirsel iptal koşulundan türetilemeyeceğini gerekçelendirmiş ve doğru özdeşlikle, doğrulanmış parametre setine bağlı düzeltilmiş analitik katsayılar üretmiştir. Harmonik etkinin deney ve kalibre benzetimdeki büyüklüğü AS3 yanıtında verilmiştir. Katkının ikinci ayağı kararlılıktır: terim akım çevriminin dışında kaldığı için marjları değiştirmez. Analog eşdeğer çevrim \SI{44}{\micro\second} gecikme marjına sahiptir. İleri besleme genliği, akım düzeyine göre ayarlanabilen bir tasarım parametresidir. Sin\textsuperscript{3} terimi modülasyon referansına eklenen tek bir ileri besleme bileşeni olduğundan mevcut akım denetleyicisi değiştirilmeden uygulanabilmektedir.

\subsection{GT-Tipi Mimari Katkısı}

GT-tipi yaklaşım, T-tipi dönüştürücü anahtarlarının evirici, PFC, izole şarj ve batarya bazlı yönetim görevlerinde yeniden kullanılmasını sağlayan bir mimari çerçeve sunmaktadır. Mimari, tezin \SI{20}{\kilo\volt\ampere}--\SI{25}{\kilo\watt} deney düzeneğinde kurulmuştur. Evirici kipi deneylerle, şarj yönü benzetim kayıtlarıyla gösterilmiştir. Mimari, ayrı PFC ve şarj dönüştürücüsü yerine aynı IGBT'leri RL1 ile yeniden görevlendirerek sürücü, koruma ve soğutma altyapısını ortak kullanma ve maliyet ile hacmi azaltma olanağı sunmaktadır.

Üç katkı birlikte değerlendirildiğinde tezin ana ekseni, T-tipi donanımın iki ölçekte yeniden kullanılması ve faz kolu ölçeğindeki değişikliğin Sin\textsuperscript{3} ileri beslemesiyle tamamlanmasıdır. Faz kolu ölçeğinde nötr kolun düşük güçteki görevi değiştirilmekte, sistem ölçeğinde dönüştürücü anahtarları şarj--deşarj mimarisi içinde yeniden görevlendirilmektedir. Sin\textsuperscript{3} terimi, faz kolu değişikliğinin harmonik bedelini akım çevriminin kararlılığını bozmadan klasik anahtarlama düzeyine indirmekte ve susturmanın verim kazancını korumaktadır.

\subsection{Fikrî Mülkiyet Çıktılarının Kanıt Konumu}

Tez çalışmasıyla ilişkili üç ulusal patent ve üç PCT başvurusu Tablo~\ref{tab:fikri-mulkiyet}'te verilmiştir. Numara, tarih, hak sahibi ve buluş sahipleri Türk Patent ve Marka Kurumu Resmî Patent Bülteni ile WIPO yayın sayfalarından okunmuştur (EK B ve EK C). Üç ulusal patentin ve üç başvurunun başlığı ``Değiştirilmiş T Tipi İki Yönlü Dönüştürücüden Beslenen Aktif Batarya Yönetim Sistemi''dir; hak sahibi Sakarya Üniversitesi Rektörlüğü, buluş sahipleri Ümit Pekparlak ve Uğur Arifoğlu'dur. Her PCT başvurusu bir ulusal başvurunun rüçhanını taşır.

\begin{table}[H]
  \centering
  \caption{Tez çalışmasıyla ilişkili fikrî mülkiyet çıktıları (kaynak: EK B ve EK C).}
  \label{tab:fikri-mulkiyet}
  \small
  \begin{tabularx}{\textwidth}{p{0.18\textwidth}p{0.24\textwidth}X}
    \toprule
    \textbf{Tür} & \textbf{Numara} & \textbf{Başvuru, veriliş ve yayın bilgisi} \\
    \midrule
    Ulusal patent & TR 2023 013377 B & Başvuru 19.10.2023; patent belgesi veriliş 21.02.2025 (Bülten 2025/2) \\
    Ulusal patent & TR 2023 013378 B & Başvuru 19.10.2023; patent belgesi veriliş 23.06.2025 (Bülten 2025/6) \\
    Ulusal patent & TR 2023 013379 B & Başvuru 19.10.2023; patent belgesi veriliş 21.02.2025 (Bülten 2025/2) \\
    PCT başvurusu & WO 2025/085037 A1 & PCT/TR2024/051205; başvuru 18.10.2024; yayın 24.04.2025; rüçhan TR 2023/013377 \\
    PCT başvurusu & WO 2025/085038 A3 & PCT/TR2024/051206; başvuru 18.10.2024; yayın 24.04.2025; arama raporu 03.07.2025; rüçhan TR 2023/013378 \\
    PCT başvurusu & WO 2025/085039 A1 & PCT/TR2024/051207; başvuru 18.10.2024; yayın 24.04.2025; rüçhan TR 2023/013379 \\
    \bottomrule
  \end{tabularx}
\end{table}

GT-tipi mimarinin temelini oluşturan değiştirilmiş T-tipi iki yönlü dönüştürücü ve aktif batarya yönetim sistemi buluşları, Türk Patent ve Marka Kurumu tarafından üç ulusal patent belgesiyle tescil edilmiştir. Buluşlar üç PCT başvurusuyla uluslararası yayın aşamasına taşınmış, WO 2025/085038 başvurusu için uluslararası arama raporu da yayımlanmıştır. Fikrî mülkiyet çıktıları, tezin teknolojik yenilik katkısını tamamlamaktadır.

\section{Sonuçların Elde Edildiği Koşullar}

Deneysel sonuçlar güvenli laboratuvar koşulunda, şebeke gerilimi varyak üzerinden yaklaşık \SI{120}{\volt} düzeyine ayarlanarak ve yaklaşık \SI{2.3}{\ampere} faz akımında elde edilmiştir. Orta (\SI{10}{\ampere}) ve anma (\SI{25}{\ampere}) akımı sonuçları S4 benzetim modeliyle üretilmiştir. S4 modeli, deneye kalibre edilen S3 modeline akım geri beslemesi eklenerek kurulmuş ve anma akımı benzetimleri analog eşdeğer akım denetimi ile büyütülmüş DA bara koşulunda yürütülmüştür.

\section{Gelecek Çalışmalar}

Öncelikli gelecek çalışma, S4 akım denetiminin sayısal gerçeklemesidir: kapasitör akımı geri beslemesiyle aktif sönümleme ya da şebeke akımı geri beslemesi tasarlanmalı, örnekleme ve hesaplama gecikmesi dahil marjlar yeniden hesaplanmalıdır. Aynı aşamada nötr susturmadan tamamlayıcı anahtarlamaya akım düzeyine bağlı kip geçişi ve Sin\textsuperscript{3} ileri besleme genliğinin akımla ayarlanması tasarlanmalı, etkinleşme eşiği, histerezis ve geçiş sürekliliği raporlanmalıdır. $L_2$ şebeke ve varyak eşdeğeri doğrudan empedans ölçümüyle, $L_1$ ise doyum eğrisiyle modellenmelidir.

İkinci aşamada düşük, orta ve güvenli biçimde erişilebilen en yüksek/nominal güç noktalarında ortak karşılaştırma yapılmalıdır. Verim, güç faktörü, THD, $I_3/I_1$, RMS akım, ters akım alanı, anahtarlama geçiş sayısı, yarı iletken sıcaklığı ve modülasyon payı en az üç tekrarla raporlanmalıdır. Akım referansı basamağı, şebeke gerilimi değişimi ve kompanzasyonun açma/kapama anı için yerleşme süresi, aşım, doyum ve koruma davranışı ayrıca ölçülmelidir. Benzetim tarafında ölçülen şebeke gerilimi harmonikleri ile şebeke eşdeğerinin $L$ ve $R$ değerleri modele eklenerek S1 faz akımı dalga biçimi yeniden değerlendirilmelidir.

Üçüncü aşamada donanımın gerçek denetim yazılımı ve örnekleme yapısıyla $L$, $S$ ve $T$ fonksiyonları anahtarlamalı modelde ve donanımda enjeksiyonla ölçülmelidir. Nötr susturmanın periyodik kip değiştirmesi harmonikler arası bağlaşımı belirginleştirirse sonlu mertebeli HTF modeli kurulmalı; kesme mertebesi ve yaklaşım hatası raporlanmalıdır. Sin\textsuperscript{3} katsayıları yalnız düşük güç THD minimumuna göre değil; temel akım, RMS kayıp, doyum, kararlılık ve geçiş sürekliliği birlikte gözetilerek seçilmeli, Alt Bölüm~\ref{sec:sin3-yeniden-tarama} taraması gerçek kod ve anahtarlamalı modelle yinelenmelidir.

GT-tipi mimari için tam ölçekli prototipten önce donanım-içinde-döngü ve düşük gerilimli güç katı doğrulaması yapılmalıdır. Ön şarj, PFC/CC--CV şarj, şebekeye deşarj, RL1 kip geçişi ve güvenli kapalı durumları; yanlış kapı komutu, sensör kaybı, aşırı akım, aşırı gerilim, sıcaklık ve haberleşme hatalarıyla sınanmalıdır. Çok sekonderli transformatörün izolasyon, kaçak endüktans ve üretilebilirlik özellikleri ile röle ve doğrultucu/boost diyotlarının kayıpları tasarımın sayısal karşılaştırmasına eklenmelidir.

Batarya çalışmaları farklı kimya, kapasite ve yaşlanma koşullarını içerecek biçimde genişletilebilir. CC--CV sınırları, SoC/SoH kestirimi, bataryalar arası akım paylaşımı ve uzun süreli termal davranış üretici güvenlik sınırları içinde değerlendirilmelidir.

Son aşamada yöntemlerin mikro şebeke, elektrikli araç şarj altyapısı ve hibrit yenilenebilir enerji sistemi gibi uygulamalara aktarılabilirliği araştırılabilir. Farklı topolojilere genişletme, T-tipi sistemde kanıtlanan mekanizma ve ölçütler üzerine kurulabilir. Yapay zekâ veya adaptif optimizasyon gibi ilave yöntemler, tezde kurulan fiziksel model ve deney zinciriyle birlikte kullanıldığında katsayı seçimini ve çalışma noktası uyarlamasını destekleyebilir.

\section{Genel Sonuç}

Tezin ulaştığı genel sonuç, T-tipi dönüştürücünün düşük güçteki nötr kol davranışının iyileştirme olanağı sunduğu ve susturmanın getirdiği verim--harmonik ödünleşiminin Sin\textsuperscript{3} ileri beslemesiyle dengelendiğidir. Deney noktasında susturma verimi 1,7 puan artırmış ve şebeke akımı THD'si \%15'ten \%47'ye çıkmıştır. Sin\textsuperscript{3} ileri beslemesi THD'yi \%14'e, yani klasik anahtarlama düzeyine indirirken verimi korumuş ve \%94,5'e taşımıştır. Uygun ileri besleme katsayıları çalışma noktasına bağlıdır. Deney özetine kalibre edilen benzetim modeli S3'te deneyin akım, verim ve THD değerlerini yeniden üretmiş, üç senaryoda verim ve THD sıralamasını korumuştur. S4 benzetiminde ileri besleme ile PI ve PIR akım denetimi, analog eşdeğer denetim ve büyütülmüş DA bara koşulunda \SI{25}{\ampere} anma akımını doyumsuz izlemiştir. Sin\textsuperscript{3} teriminin kararlılık marjlarını değiştirmediği analitik, frekans alanı ve zaman alanı kanıtlarıyla gösterilmiştir. GT-tipi yaklaşım, aynı T-tipi donanımı şebekeye deşarj, şebekeden şarj ve batarya bazlı şarj denetimi işlevlerinde kullanan tutarlı bir enerji depolama mimarisi sunmuştur. Sonuç olarak tez, T-tipi dönüştürücünün düşük güç verimini artıran, şebeke akımının harmonik kalitesini koruyan ve aynı donanımı enerji depolama sisteminin şarj ve deşarj işlevlerinde birleştiren bütünlüklü bir yaklaşım ortaya koymaktadır.
